\section{Integral distribution, reflection and modular specialization}
\label{ir:sec:scope}
The character basis above gives efficient coordinates over a splitting
field. To settle the integral question, retain the original grid symbols:
Fourier transport is exact over its cyclotomic coefficient field, but
inverse transforms introduce group-order denominators and need not be
unimodular over integral rings. The following division-free construction
proves an all-ring basis before imposing reflection. The resulting
reflection quotients can have two-torsion even though the unreflected
module is free. This develops the two overlapping integral continuations
through a single proof chain: point basis, resolution, fixed-point
homology, scalar torsion, jets and analytic identities.

The ordinary distribution and its sign cohomology are classical
\cite{dist:Kubert,intdist:KubertSign}; Ouyang's universal norm distribution
and Anderson-type resolution supply broader antecedents
\cite{dist:Ouyang}. We give the complete proofs for the specified weighted
presentation. Neither formal freeness nor a formal torsion calculation
asserts independence of evaluated periods.
\newcommand{\intdistL}[2]{\mathcal L_{#1}(#2)}
\newcommand{\intdistJ}[2]{\mathcal L^{[#1]}(#2)}

For a positive integer $q$, let
\[
 X_q=(q^{-1}\Z)/\Z,\qquad P(q)=\{p:p\text{ is prime and }p\mid q\},
 \qquad R_q=\Z[A_p:p\in P(q)].
\]
We write $\operatorname{den}(x)$ for the exact order of $x$ in $\Q/\Z$, and set
$\operatorname{den}(0)=1$. A symbol $e_x$ is not itself a special value. It is a generator
in a free module. Under polylogarithm evaluation, $e_0$ will become
$\Li_s(1)=\zeta(s)$. Under Hurwitz evaluation, it becomes
$\zeta(s,1)$, not the undefined expression $\zeta(s,0)$.

For $d\mid q$, put $A_d=\prod_{p\mid d}A_p^{v_p(d)}$.
Weights on composite divisors are thus multiplicative; an unrelated
variable for every divisor would define a different module. The universal
module and its relation submodule are
\begin{align}
 V_q&=\bigoplus_{x\in X_q}R_qe_x,\\
 r_{p,x}&=\sum_{\substack{y\in X_q\\py=x}}e_y-A_pe_x
       \qquad(p\mid q,\ x\in X_{q/p}),\label{ir:eq:prime-row}\\
 \mathcal R_q&=\langle r_{p,x}\rangle_{R_q},\qquad
 \mathcal D_q=V_q/\mathcal R_q.\label{ir:eq:module}
\end{align}
Every sum in \eqref{ir:eq:prime-row} has exactly $p$ terms. For a
commutative ring $S$ and weights $a_p\in S$, write $\mathcal D_q(S,a)$ for the
same presentation after $A_p\mapsto a_p$.

\begin{lemma}[Prime rows suffice]\label{ir:lem:divisor}
The relation $r_{d,x}=\sum_{dy=x}e_y-A_de_x$, for $d\mid q$ and
$x\in X_{q/d}$, follows from the prime rows without division.
\end{lemma}
\begin{proof}
For $ab\mid q$ and $x\in X_{q/(ab)}$, direct expansion gives
\[
 r_{ab,x}=\sum_{ay=x}r_{b,y}+A_b r_{a,x}.
\]
Here all $y$ belong to $X_{q/b}$, so every displayed prime or divisor
row is legitimate. The intermediate symbols cancel. Induction on the
number of prime factors of $d$, counted with multiplicity, proves the
claim.
\end{proof}

This elementary observation already prevents a common ambiguity:
``distribution rank'' is meaningful only after the row vocabulary,
endpoint convention, coefficient ring, and divisor-weight convention
have all been fixed.
\subsection{A division-free integral basis and a unit relation minor}
\label{ir:sec:integral}

\subsubsection{Choosing the surviving grid points}
Write $U_m=(\Z/m\Z)^\times$ and use a one-element set $U_1$.
For every prime $p$ and exponent $a\ge1$, choose a distinguished lift
in each fiber of $U_{p^a}\to U_{p^{a-1}}$. At $a=1$, this means choosing
one member of $U_p$. Let $E_{p,a}$ be the remaining members. Thus
\begin{equation}\label{ir:eq:local-count}
 |E_{p,a}|=\varphi(p^a)-\varphi(p^{a-1}),\qquad \varphi(1)=1.
\end{equation}
The following fixed choice makes the construction executable:
\begin{align}
 E_{p,1}&=\{2,3,\ldots,p-1\},\label{ir:eq:local-choice}\\
 E_{p,a}&=\{b+jp^{a-1}:1\le b<p^{a-1},\ p\nmid b,
                         \ 1\le j\le p-1\}\quad(a\ge2).\nonumber
\end{align}
For $p=2$, the first set is empty. This is intentional.

Every point $x$ of exact denominator $m$ has a unique additive primary
expansion
\begin{equation}\label{ir:eq:primary-coordinates}
 x=\sum_{p^a\parallel m}\frac{b_p}{p^a}\pmod\Z,
 \qquad b_p\in U_{p^a}.
\end{equation}
For $x=n/m$ in lowest terms,
$b_p\equiv n(m/p^a)^{-1}\pmod{p^a}$. Define $\mathcal B_q$ to contain $0$
and all points of denominator $m\mid q$ for which
$b_p\in E_{p,a}$ at every nonzero primary component. We call a primary
component \emph{bad} when its residue is the distinguished, discarded
lift. This adjective refers only to a basis choice.

The cardinality telescopes prime by prime:
\begin{equation}\label{ir:eq:basis-count}
 |\mathcal B_q|=
 \prod_{p^e\parallel q}\left(1+
   \sum_{a=1}^{e}\bigl(\varphi(p^a)-\varphi(p^{a-1})\bigr)\right)
 =\varphi(q).
\end{equation}
The point $0$ must be retained for this count and for the normal forms.

\begin{theorem}[Integral normal form]\label{ir:thm:integral}
The images of $e_b$, $b\in\mathcal B_q$, form an $R_q$-basis of $\mathcal D_q$.
Consequently, for every commutative ring $S$ and every choice of weights
$a_p\in S$, the same points form an $S$-basis of $\mathcal D_q(S,a)$.
In particular this module is free of rank $\varphi(q)$ even in positive
characteristic, at zero weights, and over nonreduced rings.
\end{theorem}
The spanning algorithm is immediate and is proved next. Its independence
is established in Section~\ref{ir:sec:resolution}, without dividing by
any group order or weight.

\subsubsection{The terminating straightening rule}
Let $y\notin\mathcal B_q$, let $m=\operatorname{den}(y)$, and choose any bad prime $p$ for
its primary coordinates. The row with target $py$ gives
\begin{equation}\label{ir:eq:straighten}
 e_y=A_pe_{py}-\sum_{\substack{pz=py\\z\ne y}}e_z
       \quad\text{in }\mathcal D_q.
\end{equation}
All terms on the right have denominator dividing $m$.
At exact denominator $m$, the $p$-primary components of the other
solutions are the allowed lifts in the same fiber; all other primary
components are unchanged. When the $p$-component has order $p$, there
is also one solution with zero $p$-component, hence lower denominator.
The point $py$ itself has denominator $m/p$.

Order points first by their positive denominator and then by the number
of bad primary components. Every right-hand term of
\eqref{ir:eq:straighten} is strictly earlier in this order. Recursion
therefore terminates, and proves that $\mathcal B_q$ spans. The implementation
always chooses the smallest bad prime, but the theorem permits any
choice. Independence will imply that all choices yield the same
normal form.

\begin{proposition}[Support and polynomial degree]\label{ir:prop:support}
Write the normal form of $e_y$, with $m=\operatorname{den}(y)$, as
\[
 \operatorname{NF}(e_y)=\sum_{b\in\mathcal B_q}c_{y,b}(\mathbf A)e_b.
\]
A nonzero coefficient has $\operatorname{den}(b)=d\mid m$, and
\begin{equation}\label{ir:eq:degree-bound}
 \deg_{A_p}c_{y,b}\le v_p(m/d)\qquad(p\mid q).
\end{equation}
All coefficients belong to $\Z[\mathbf A]$.
\end{proposition}
\begin{proof}
The recursion never increases a primary denominator. Multiplication by
$A_p$ occurs only in the term $e_{py}$, where that denominator has
lost one power of $p$. Every path from denominator $m$ to denominator
$d$ can therefore contain at most $v_p(m/d)$ such multiplications.
The other coefficients in the recursion are $-1$. Summing the finitely
many paths proves the assertions, including possible cancellations.
\end{proof}

There is no assertion here that the coefficient of $e_0$ factors as a
product of $A_p-1$. At level $15$ the coefficient in one normal form is
$1-A_3A_5$; this distinction will matter at the spectral pole.

\subsubsection{An explicit maximal minor equal to one}
For each nonbasis point $y$, select the single row $r_{p,py}$ used in
its first straightening step. Select the nonbasis points as columns,
and order rows by their selected pivot. The restricted matrix is
triangular, with every diagonal entry equal to $1$: all other nonbasis
columns are earlier in the terminating order.

The selected rows are distinct. Indeed a selected row has its chosen
point as its unique latest nonbasis point, with coefficient $1$. Two
identical rows could not have different such latest points.
We have proved the following stronger integrality statement.

\begin{theorem}[Universal unit minor and relation basis]
\label{ir:thm:minor}
The prime-relation matrix contains an explicitly selected
$(q-\varphi(q))\times(q-\varphi(q))$ minor equal to $1$ in $R_q$.
The selected rows form an $R_q$-basis of the full relation module
$\mathcal R_q$. After specializing the weights to integers, all nonzero Smith
invariants of the complete prime-relation matrix are $1$; there are
exactly $q-\varphi(q)$ of them.
\end{theorem}
\begin{proof}
Triangularity proves the minor statement. Selected-row elimination
reduces every vector to a combination of basis-point symbols. By
Theorem~\ref{ir:thm:integral}, a vector in $\mathcal R_q$ has zero such
remainder, so every raw row is a combination of the selected ones.
The minor proves their independence. The same argument specializes
without division. Equivalently, after an integer specialization the
image lattice is saturated and has the indicated rank; the unit minor
also proves directly that all Smith invariants are units.
\end{proof}
The proof order is not circular: the resolution below proves
Theorem~\ref{ir:thm:integral} independently of the assertion that the
selected rows span all relations.

\subsubsection{Certificates rather than opaque elimination}
The straightener can retain the exact row combination at every step.
If $C_y$ denotes a certificate for $e_y-\operatorname{NF}(e_y)$, then
\begin{equation}\label{ir:eq:certificate-recursion}
 C_y=r_{p,py}+A_pC_{py}-\sum_{\substack{pz=py\\z\ne y}}C_z,
 \qquad C_b=0\ (b\in\mathcal B_q).
\end{equation}
A verifier need not trust the recursion. It can expand the listed raw
rows in the original $q$-generator free module, multiply their explicit
integer polynomials, and compare coefficients. The accompanying
verifier does exactly this for the three frozen identities.

\begin{example}[Small levels]
The specified basis choices give
\[
 \mathcal B_4=\{0,3/4\},\qquad
 \mathcal B_{12}=\{0,5/12,2/3,3/4\},
\]
\[
 \mathcal B_{15}=\{0,1/15,4/15,2/5,7/15,3/5,2/3,4/5\}.
\]
Also $\mathcal B_{30}=\mathcal B_{15}$ as sets of rational points. An exact primary
component of order two has no allowed residue, which explains this
last equality. It does not make the weight $A_2$ irrelevant to all
normal-form coefficients.
\end{example}

The procedure terminates with polynomial coefficients, but we do not
assert a polynomial-time bound in the bit length of $q$. There are
$q$ raw grid points to begin with, and certificate coefficients can
grow. The finite support and degree bounds are rigorous; sharper
complexity and coefficient-height estimates remain separate questions.

\subsubsection{A uniform coefficient-height bound}
For $q>1$, let $s=|P(q)|$, let $P_{\max}$ be its largest prime divisor,
and put $H(x)=\sum_p v_p(\operatorname{den}(x))$,
$b(x)=\#\{\text{bad primary components of }x\}$ and
$W(x)=(s+1)H(x)+b(x)$. Summing absolute values over all normal-form
coordinates and all polynomial monomials gives
\[
 \|\operatorname{NF}(e_x)\|_1\le P_{\max}^{W(x)}.
\]
Indeed every straightening edge lowers $W$ by at least one: at equal
denominator a bad component is removed, while a decrease in $H$ absorbs
any possible increase of at most $s$ in $b$. There are at most
$P_{\max}$ terms in a step, and multiplication by $A_p$ preserves the
coefficient norm. Induction proves the bound. The sharper coordinate
multidegrees remain those in \eqref{ir:eq:degree-bound}. The estimate is
a uniform height bound, not an optimal complexity assertion.

\subsubsection{The complete level-twelve normal-form table}
\label{ir:tab:complete12}
Let $E_a=[e_{a/12}]$, $\alpha=A_2$, $\beta=A_3$.
The retained basis is $E_0,E_5,E_8,E_9$. Each row gives the four
coefficients of $E_a$ in that basis.
\begin{center}
\renewcommand{\arraystretch}{1.28}
\begin{tabular}{@{}c r r r r@{}}
\toprule
$a$&$E_0$&$E_5$&$E_8$&$E_9$\\
\midrule
0&$1$&$0$&$0$&$0$\\
1&$\alpha\beta(\alpha-1)$&$-1$&$0$&$-\beta-1$\\
2&$\alpha(\beta-1)$&$0$&$-\alpha-1$&$0$\\
3&$\alpha(\alpha-1)$&$0$&$0$&$-1$\\
4&$\beta-1$&$0$&$-1$&$0$\\
5&$0$&$1$&$0$&$0$\\
6&$\alpha-1$&$0$&$0$&$0$\\
7&$\alpha(\beta-\alpha)$&$1$&$-\alpha(\alpha+1)$&$\beta+1$\\
8&$0$&$0$&$1$&$0$\\
9&$0$&$0$&$0$&$1$\\
10&$1-\beta$&$0$&$\alpha+1$&$0$\\
11&$\alpha(1-\beta)$&$-1$&$\alpha(\alpha+1)$&$0$\\
\bottomrule
\end{tabular}
\end{center}
The companion file \texttt{normal\_forms\_q12.json} contains the
unfactored integer monomial data. At arbitrary complex order, substitute
$\alpha=2^{1-s}$, $\beta=3^{1-s}$ and
$E_a=\Li_s(e^{2\pi ia/12})$. At formal integral jet weights, substitute
polynomials modulo $u^L$ instead. The same table supports both uses.

\subsection{A finite weighted resolution and a proof of independence}
\label{ir:sec:resolution}

\subsubsection{The complex}
Order the primes dividing $q$. For $S\subseteq P(q)$ put
$p_S=\prod_{p\in S}p$, and define
\begin{equation}\label{ir:eq:chain-groups}
 C_j=\bigoplus_{\substack{S\subseteq P(q)\\|S|=j}}
       R_q[X_{q/p_S}].
\end{equation}
A generator is written $[S,x]$. Define
\begin{align}
 d[S,x]=\sum_{p\in S}(-1)^{\#\{\ell\in S:\ell<p\}}
 \left(\sum_{py=x}[S\setminus\{p\},y]
       -A_p[S\setminus\{p\},x]\right).
 \label{ir:eq:chain-differential}
\end{align}
The roots in each inner sum belong to $X_{q/p_{S\setminus\{p\}}}$.
The augmentation $C_0\to\mathcal D_q$ is the defining quotient.

\begin{lemma}\label{ir:lem:complex}
The operators in distinct prime directions commute before the exterior
signs are applied. Hence $d^2=0$, and the image of $C_1\to C_0$ is
exactly $\mathcal R_q$.
\end{lemma}
\begin{proof}
Both successive root sums in prime directions $p$ and $\ell$ enumerate
once all solutions of $p\ell z=x$. The two terms consisting of one
root sum and one multiplication by a weight agree in either order,
and so do the terms with coefficient $A_pA_\ell$. All intermediate
points lie in the required grids. The exterior signs make each pair
cancel in $d^2$. The assertion about $C_1$ is its definition.
\end{proof}

\subsubsection{A filtration that removes the weights}
Associate to a generator the integer
\begin{equation}\label{ir:eq:augmented-conductor}
 \nu(S,x)=p_S\operatorname{den}(x),
\end{equation}
which divides $q$. Filter the complex by increasing $\nu$.
If $p\in S$, a root of $py=x$ has denominator either
$\operatorname{den}(x)$ or $p\operatorname{den}(x)$. Thus the differential does not increase
$\nu$. Its top-filtration part retains exactly the roots of denominator
$p\operatorname{den}(x)$; the term with coefficient $A_p$ lies strictly lower.
The associated graded differential is consequently independent of all
weights.

Fix a divisor $m\mid q$. In its associated graded summand, a generator
with subset $S$ has exact denominator $m/p_S$, so necessarily
$S\subseteq P(m)$. For $p^a\parallel m$, the local coordinate belongs
to $U_{p^{a-1}}$ when $p\in S$, and to $U_{p^a}$ otherwise. At exponent
zero the unique local coordinate is zero, represented by the one-element
set $U_1$.

There is a harmless CRT twist: division by a prime changes the other
primary coordinates by a unit. We spell out its removal because omitting
it would leave a gap in a claimed tensor-product argument. If $x_\ell$
is the $\ell$-primary component of $x$, replace it in the $S$-summand by
\begin{equation}\label{ir:eq:gauge}
 z_\ell=\left(\prod_{p\in S\setminus\{\ell\}}p\right)^{-1}x_\ell.
\end{equation}
The inverse exists on the $\ell$-primary group. Under a differential
edge removing $p$, the coordinates $z_\ell$ for $\ell\ne p$ are
unchanged: the factor $p^{-1}$ in a root coordinate cancels the removal
of $p$ from the product. The $p$-coordinate is carried to all its lifts
under multiplication by $p$.

It follows that the associated graded summand at $m$ is the tensor
product, with its usual exterior signs, of the two-term complexes
\begin{equation}\label{ir:eq:local-complex}
 R_q[U_{p^{a-1}}]\xrightarrow{\,N_{p,a}\,}R_q[U_{p^a}],
 \qquad
 N_{p,a}(b)=\sum_{\substack{u\in U_{p^a}\\u\mapsto b}}u,
 \quad p^a\parallel m.
\end{equation}
For $a=1$, the image of the single generator is the sum of all $p-1$
nonzero residues. For $a\ge2$, each fiber has $p$ elements.

\begin{lemma}[Local splitting]\label{ir:lem:local-splitting}
Each $N_{p,a}$ is a split injection. Its cokernel is free with basis the
images of $E_{p,a}$.
\end{lemma}
\begin{proof}
Different fibers have disjoint supports. In one fiber, the sum of its
basis vectors has a coefficient equal to $1$ on the distinguished lift.
Replace that one vector by the fiber sum; this is a unimodular basis
change. The remaining vectors are exactly the allowed lifts. Doing this
in all fibers proves both claims over any coefficient ring.
\end{proof}

\begin{theorem}[Exact weighted resolution]\label{ir:thm:resolution}
The augmented complex
\[
 0\longrightarrow C_{|P(q)|}\longrightarrow\cdots
 \longrightarrow C_1\longrightarrow C_0\longrightarrow\mathcal D_q
 \longrightarrow0
\]
is exact. The images of $\mathcal B_q$ are a basis of $\mathcal D_q$, proving
Theorem~\ref{ir:thm:integral}.
\end{theorem}
\begin{proof}
By Lemma~\ref{ir:lem:local-splitting}, each local two-term complex is
a direct sum of a contractible identity complex and its free cokernel
in degree zero. Its finite tensor product is therefore acyclic in
positive degrees. Its degree-zero homology has the tensor-product
basis of the allowed local residues. In the graded summand at $m$,
these are exactly the basis points of denominator $m$.

Now lift through the finite filtration. If a positive-degree cycle
has highest nonzero filtration $m$, its highest component is a
boundary in the associated graded complex. Subtract a lift of that
boundary to obtain a cycle of lower filtration. Repetition makes the
cycle a boundary in the original complex. This proves positive-degree
exactness without any convergence issue.

For degree zero, suppose a linear combination of allowed points is
$d w$ with $w\in C_1$, and let $M$ be the highest filtration present
in that linear combination. If $w$ has higher filtration, its leading
component is a cycle in degree one of the associated graded complex.
Positive-degree exactness there allows us to subtract a lifted boundary
$d v$, $v\in C_2$, from $w$ and lower its filtration without changing
$d w$. Repetition gives a preimage of filtration at most $M$. Its
leading boundary would now be a relation among the allowed points in
the graded summand at $M$. Those points form a free basis of its
zero-degree homology, so all their coefficients must vanish. Descending
$M$ proves independence. Spanning was proved by
\eqref{ir:eq:straighten}. Thus the claimed points form a basis of
$H_0(C_\bullet)=\mathcal D_q$.
\end{proof}

\subsubsection{Base change and the rank identity}
The proof works verbatim over an arbitrary commutative ring with
arbitrary weights. Alternatively, the augmented complex over $R_q$
is split exact as an underlying $R_q$-module complex: its degree-zero
homology is free, and successive surjections onto projective kernels
split. Tensoring with any $R_q$-algebra therefore preserves exactness.
No flatness hypothesis on the target algebra is needed.

Taking the alternating ranks of the chain groups gives a useful check:
\begin{equation}\label{ir:eq:euler-rank}
 \rank\mathcal D_q
 =\sum_{S\subseteq P(q)}(-1)^{|S|}\frac{q}{p_S}
 =q\prod_{p\mid q}(1-p^{-1})=\varphi(q).
\end{equation}
This rank computation alone would not prove integral freeness. The
local unimodular splittings and the lifted basis are the missing
integral information.

\begin{remark}[Relation to character normal forms]
Over a characteristic-zero splitting field, both the conductor-character
basis in the source and the present raw-point basis describe the same
free module. Their transition matrix is a legitimate field-valued
change of coordinates. Its construction can involve divisions by
character-group orders. The raw-point basis is designed to avoid those
divisions, not to replace the arithmetic convenience of character
coordinates where those coordinates are available.
\end{remark}

\subsection{Reflection and the exact integral torsion formula}
\label{ir:sec:reflection}

\subsubsection{Reflection is a separate quotient operation}
The map $c(e_x)=e_{-x}$ preserves the relations and satisfies $c^2=1$.
It acts on the entire resolution by $c[S,x]=[S,-x]$.
For integer weights $a_p$, let $D=\mathcal D_q(\Z,a)$ and define
\[
 D^+=D/(c-1)D,\qquad D^-=D/(c+1)D.
\]
These are \emph{coinvariants with a prescribed sign}, not the submodules
of invariant and anti-invariant elements. Over a field of characteristic
not two, the corresponding eigenspaces and coinvariants agree. Over
$\Z$ this distinction is essential.

For any $q$, put
\begin{equation}\label{ir:eq:r-h}
 r(q)=\#\{p>2:p\mid q\}+\mathbf 1_{4\mid q}.
\end{equation}
For $q>2$, we have $r(q)\ge1$ and set
$n(q)=\varphi(q)/2$, $h(q)=2^{r(q)-1}$.

\begin{theorem}[Weighted integral reflection theorem]
\label{ir:thm:integer-reflection}
Let $q>2$ and specialize every $A_p$ to an integer $a_p$. For each
$\eps\in\{1,-1\}$,
\begin{equation}\label{ir:eq:integer-reflection}
 D/(c-\eps)D\simeq\Z^{n(q)}\oplus(\Z/2\Z)^{t(q,a)},
\end{equation}
where
\begin{equation}\label{ir:eq:torsion-count}
 t(q,a)=
 \begin{cases}
  h(q),&a_p\text{ is odd for every odd prime }p\mid q,\\
  0,&\text{some odd prime }p\mid q\text{ has }a_p\text{ even}.
 \end{cases}
\end{equation}
In particular the answer is independent of the integer $a_2$.
At $q=1$ or $q=2$, the separate answers are $D^+=\Z$ and
$D^-=\Z/2\Z$.
\end{theorem}

The absence of odd torsion is an elementary property of an involution
on a free lattice. The exact multiplicity is the substantive calculation.
We prove it through the fixed points of the resolution.

\subsubsection{Two elementary facts about involutions}
For a module $M$ with an involution $c$, define its two periodic
cohomology groups by
\begin{equation}\label{ir:eq:tate-definition}
 \widehat H^0(M)=\ker(c-1)/(1+c)M,\qquad
 \widehat H^1(M)=\ker(1+c)/(c-1)M,
\end{equation}
and repeat with period two. These are the cohomology groups of the
complex whose differentials alternate between $c-1$ and $1+c$,
with $c-1$ leaving even degree. This definition, rather than any
unmentioned cohomology convention, fixes all signs and parity labels.

\begin{lemma}[Torsion in sign coinvariants]\label{ir:lem:torsion-tate}
If $M$ is a finite-rank free abelian group, every torsion class in
$M/(c-\eps)M$ is killed by $2$. Its torsion subgroup is
$\widehat H^1(M)$ for $\eps=1$, and $\widehat H^0(M)$ for $\eps=-1$.
\end{lemma}
\begin{proof}
Suppose $k v=(c-\eps)w$ with $k\ne0$. Applying $c+\eps$ gives
$k(c+\eps)v=0$, hence $(c+\eps)v=0$ because $M$ is torsion-free.
Then $(c-\eps)v=-2\eps v$, so the class of $v$ is killed by $2$.
Conversely a vector in $\ker(c+\eps)$ represents a class killed by
$2$. Taking its quotient by $(c-\eps)M$ gives exactly the two groups
in \eqref{ir:eq:tate-definition}.
\end{proof}

\begin{lemma}[Balanced free ranks]\label{ir:lem:balanced}
For every field of characteristic different from two, and arbitrary
weights in that field, the two reflection coinvariants of $\mathcal D_q$
have dimension $\varphi(q)/2$ when $q>2$.
\end{lemma}
\begin{proof}
On a permutation module with $u$ basis points and $f$ fixed points,
the two eigenspace dimensions are the integers $(u+f)/2$ and
$(u-f)/2$. The sign-eigenspace functors are exact when $2$ is
invertible, so apply this count to the resolution. The alternating
sum of the $u$'s is $\varphi(q)$ by \eqref{ir:eq:euler-rank}.
The fixed-point count for $X_m$ is one if $m$ is odd and two if
$m$ is even. Thus its alternating sum is
\[
 \sum_{S\subseteq P(q)}(-1)^{|S|}\gcd(2,q/p_S)=0
 \qquad(q>2).
\]
For odd $q$, this is $(1-1)^{|P(q)|}$. If $v_2(q)=1$, summing
first over whether $2\in S$ leaves $(2-1)(1-1)^k$, where
$k$ is the number of odd prime divisors and is positive. If
$v_2(q)\ge2$, that first sum is $2-2$. Hence both dimensions
are $\varphi(q)/2$.

This uses integer dimension counts in the chain modules, not merely
a trace equality in the coefficient field. The latter alone would
not determine dimensions in positive characteristic.
\end{proof}

\subsubsection{Removing every nonfixed orbit from periodic cohomology}
This step explains why the calculation is small despite the $q$
original grid points. Let $C^{\mathrm{mov}}_j$ be the span of the
nonfixed points in each summand of $C_j$. It is a subcomplex:
if $x\ne-x$ and $py=x$, then $y$ cannot be fixed; the term
$-a_p[S\setminus\{p\},x]$ is also nonfixed. Each term of this
subcomplex is a direct sum of copies of the regular permutation
module $S[\Z/2\Z]$ over the chosen coefficient ring $S$.

The periodic complex of a regular two-point orbit is exact over
\emph{every} commutative ring. Indeed the kernel of the difference
map is the diagonal, which is the image of the sum map; the kernel
of the sum map consists of $(u,-u)$, which is the image of the
difference map. This assertion requires neither division by $2$
nor absence of $2$-torsion.

Let $K_\bullet=C_\bullet/C^{\mathrm{mov}}_\bullet$. Its generators
are only the fixed symbols $[S,0]$ and, when available, $[S,1/2]$.
The action on $K_\bullet$ is trivial. Forming the double complex of
the periodic cohomology complexes gives
\begin{equation}\label{ir:eq:tate-comparison}
 \widehat H^\bullet(\mathcal D_q(S,a))\simeq
 H^\bullet\operatorname{Tot}\bigl(T(K_\bullet)\bigr).
\end{equation}
Here $T$ denotes the alternating difference/sum complex just defined,
and the chain degree of $K_j$ contributes $-j$ to total degree.

For completeness, this comparison can be justified without assuming a
spectral sequence degenerates. The augmented resolution is exact after
base change, so taking horizontal homology in the double complex
replaces $C_\bullet$ by $D$ in degree zero. Removing the moving
subcomplex does not change total cohomology because each of its
vertical columns is exact. In both statements the horizontal width
is bounded by $|P(q)|$, so in each total degree only finitely many
columns occur; successive elimination along these columns justifies
the two comparisons. In particular, no convergence of an infinite
horizontal filtration is being used.

\subsubsection{The fixed-point complex is a Koszul complex}
For a sequence $b_1,\ldots,b_r$ in a ring $S$, let
$K_S(b_1,\ldots,b_r)$ denote the tensor product of the two-term
complexes $S\xrightarrow{b_i}S$, with source in chain degree one.
This is the Koszul complex; its definition is all that is needed below.

For an odd prime $p$, the fixed-point part of its distribution map
is multiplication by $1-a_p$ on each available fixed symbol. Of the
$p$ roots of a fixed point, exactly one is fixed, and is the original
point. The other roots disappear in the quotient by moving orbits.
The factor at $2$ has two different forms.

When $v_2(q)=1$, its source has only $0$ and its target has $0,1/2$.
Its column is
\begin{equation}\label{ir:eq:two-column}
 \begin{pmatrix}1-a_2\\1\end{pmatrix}.
\end{equation}
The unit entry splits off an identity complex, leaving one free
copy of $S$ in degree zero. Therefore this prime contributes no
Koszul generator to the fixed-point complex.

When $v_2(q)\ge2$, both grids have two fixed points. In the order
$(0,1/2)$, the matrix is
\begin{equation}\label{ir:eq:two-matrix}
 B(a_2)=\begin{pmatrix}1-a_2&0\\1&-a_2\end{pmatrix}.
\end{equation}
The two roots of $1/2$ are $1/4,3/4$, which are moving points; this
explains the second column. Unimodular row and column operations
transform this matrix to $\operatorname{diag}(1,a_2(a_2-1))$, up to
a unit sign. These operations are valid over $\Z[a_2]$ and over
every specialization, because they pivot on the displayed $1$.
Odd-prime maps are scalar, so the operations commute with all other
directions of the complex.

We have proved the following ring-uniform reduction.
\begin{proposition}[Fixed-point reduction]\label{ir:prop:fixed-koszul}
Over any commutative ring, $K_\bullet$ is chain-homotopy equivalent to
\begin{equation}\label{ir:eq:active-koszul}
 K_S(\boldsymbol\delta),\qquad
 \boldsymbol\delta=
 \bigl(a_p-1:\ p>2,\ p\mid q\bigr)
 \ \text{together with }a_2(a_2-1)\text{ if }4\mid q.
\end{equation}
It has $r(q)$ active scalar entries. At $q=1,2$, this is the empty
Koszul complex $S$ in degree zero.
\end{proposition}

\subsubsection{Completing the integer calculation}
For $S=\Z$, the periodic complex on a trivial module alternates
between $0$ and multiplication by $2$. It separates into adjacent
two-term blocks. Tensoring one such block with the free complex
$K_\Z(\boldsymbol\delta)$ computes
$K_\Z(\boldsymbol\delta)\otimes\mathbb F_2$: multiplication by $2$ is a
free resolution of reduction modulo two. Consequently
\begin{equation}\label{ir:eq:integer-tate-koszul}
 \widehat H^\eps(D)\simeq
 \bigoplus_{\substack{0\le j\le r(q)\\j\equiv\eps\ (2)}}
 H_j\bigl(K_{\mathbb F_2}(\overline{\boldsymbol\delta})\bigr),
 \qquad\eps=0,1.
\end{equation}
Only an isomorphism of abelian groups is asserted here; no extra
multiplicative structure is required.

If some odd $a_p$ is even, its active scalar reduces to $1$, so its
two-term complex is contractible and all the homology vanishes.
Otherwise every odd-prime scalar is zero modulo two. The remaining
scalar $a_2(a_2-1)$, when present, is also zero modulo two for every
integer $a_2$. All differentials in the reduced active complex
therefore vanish, and
\[
 \dim_{\mathbb F_2}H_j=\binom{r(q)}j.
\]
For $r(q)\ge1$, the sums over even and odd $j$ are both
$2^{r(q)-1}$. Lemmas~\ref{ir:lem:torsion-tate} and
\ref{ir:lem:balanced} now prove
Theorem~\ref{ir:thm:integer-reflection}. When $r(q)=0$, the only
homology is in degree zero, giving the stated distinct answers for
$q=1,2$.

\subsubsection{Examples and the classical specialization}
At level four, put $a=a_2$. The basis is $(e_0,e_{3/4})$, and
\[
 e_{1/2}=(a-1)e_0,\qquad e_{1/4}=a(a-1)e_0-e_{3/4},
 \qquad
 c=\begin{pmatrix}1&a(a-1)\\0&-1\end{pmatrix}.
\]
For integer $a$, the nonzero Smith factor of either $c-1$ or $c+1$
is $2$, because $a(a-1)$ is even. This two-dimensional example
exhibits both the unavoidable parity and the irrelevance of the
integer value of $a_2$.

For all odd-prime weights odd, the following values result:
\[
\begin{array}{c|rrrrrr}
 q&4&12&15&30&60&105\\\hline
 n(q)&1&2&4&4&8&24\\
 r(q)&1&2&2&2&3&3\\
 h(q)&1&2&2&2&4&4
\end{array}
\]
For example, either sign at level $60$ gives
$\Z^8\oplus(\Z/2)^4$ in the exceptional parity case, but only
$\Z^8$ as soon as $a_3$ or $a_5$ is even.

When all $a_p=1$, the count agrees with the classical ordinary
sign-cohomology theorem of Kubert \cite{intdist:KubertSign}; twice an odd
level has the same effective prime count as that odd level. The
weighted parity criterion, the fixed-point matrices, and the later
jet calculation are the present explicit continuation of the
manuscript's integral question.

\subsection{Characteristic two, the universal support, and all jet lengths}
\label{ir:sec:modular}

\subsubsection{A characteristic-two homology theorem over arbitrary algebras}
Assume now that the coefficient ring $S$ has characteristic two.
Write $N=c-1=c+1$. Then $N^2=0$, and both periodic cohomology
groups are the same module $\ker N/\operatorname{im} N$. On a fixed-point module
both periodic differentials vanish. The comparison
\eqref{ir:eq:tate-comparison} therefore has a particularly simple
consequence.

\begin{theorem}[Koszul model in characteristic two]
\label{ir:thm:char2-koszul}
For every commutative characteristic-two ring $S$ and every choice
of its weights,
\begin{equation}\label{ir:eq:char2-homology}
 \frac{\ker(N:\mathcal D_q(S,a)\to\mathcal D_q(S,a))}{\operatorname{im} N}
 \simeq\bigoplus_{j=0}^{r(q)}H_j(K_S(\boldsymbol\delta)),
\end{equation}
where the active sequence is \eqref{ir:eq:active-koszul}.
\end{theorem}
\begin{proof}
Remove the moving-orbit subcomplex as in
Section~\ref{ir:sec:reflection}. On its fixed quotient the vertical
periodic differential is zero. Total cohomology is therefore the
direct sum of the horizontal homology groups, one from each of the
finitely many chain degrees. Apply the chain-homotopy reduction of
Proposition~\ref{ir:prop:fixed-koszul}.
\end{proof}

\subsubsection{The universal obstruction is a reduced complete intersection}
Let $R_2=\mathbb F_2[A_p:p\mid q]$ and let $I_q\subseteq R_2$ be the ideal
\begin{equation}\label{ir:eq:universal-ideal}
 I_q=(A_p-1:\ p>2,\ p\mid q)
       +(A_2(A_2-1):\ 4\mid q),
\end{equation}
where the second summand is omitted unless indicated. Its generators
form a regular sequence. Indeed the odd-prime entries are independent
linear variables shifted by one; after removing them, the possible
last entry is a nonzero monic polynomial in the still-free variable
$A_2$. At each step multiplication by that entry is injective in the
quotient by the earlier ones.

The Koszul complex on a regular sequence resolves its quotient. This
can be seen inductively here without a general regular-sequence
criterion: tensor a free resolution of the previous quotient with
$[R_2\xrightarrow{\delta}R_2]$. Its homology consists of the kernel
and cokernel of multiplication by $\delta$ on that quotient. The
kernel is zero, and the cokernel is the next quotient.

\begin{corollary}[Universal characteristic-two support]
\label{ir:cor:universal-support}
Over $R_2$, the reflection homology is a cyclic module:
\begin{equation}\label{ir:eq:universal-homology}
 \ker N/\operatorname{im} N\simeq R_2/I_q.
\end{equation}
Its zeroth Fitting ideal is exactly $I_q$. The closed locus defined
by $I_q$ is reduced. For odd $q>1$ it is the single point with every
weight equal to one. For $4\mid q$ it has two points, with all
odd-prime weights equal to one and $A_2=0$ or $1$. For $q$ twice an
odd number, $A_2$ is free and the locus is an affine line.
\end{corollary}
\begin{proof}
The positive Koszul homology vanishes by the preceding argument;
apply Theorem~\ref{ir:thm:char2-koszul}. A cyclic module $R_2/I_q$
has zeroth Fitting ideal $I_q$ by its one-generator presentation.
After the odd linear equations, the only possible further equation
is $A_2(A_2-1)$, whose two factors are relatively prime. Its quotient
is a product of two fields, or of two copies of the remaining
coefficient ring if one extends scalars. The stated descriptions
and reducedness follow.
\end{proof}

The cyclic universal answer does \emph{not} imply that every fiber
of the reflection homology is one-dimensional. Homology need not
commute with a nonflat specialization. At an exceptional field-valued
point the differentials of the specialized Koszul complex all vanish,
so its homology has total dimension $2^{r(q)}$. This is the precise
source of the larger fiber defect.

There is also an integral universal formulation. With the periodic
convention \eqref{ir:eq:tate-definition}, the same fixed-complex
calculation, followed by reduction modulo two, gives
\begin{equation}\label{ir:eq:universal-integral-tate}
 \widehat H^{2j}(\mathcal D_q)\simeq R_q/(2,\boldsymbol\delta),\qquad
 \widehat H^{2j+1}(\mathcal D_q)=0.
\end{equation}
Here $\mathcal D_q$ is still over the polynomial ring, not an integer weight
fiber. The regularity argument is in $R_q/2R_q$. Odd periodic
cohomology appearing after integer specialization is another
nonflat-base-change effect, not a contradiction to
\eqref{ir:eq:universal-integral-tate}.

\subsubsection{The exact rank locus over every field}
For a field $k$ of characteristic two, a nonzero active scalar is a
unit and makes the Koszul complex contractible. If all active scalars
vanish, its differentials vanish and its total homology dimension is
$2^{r(q)}$. To pass from homology to the desired coinvariants, note
that an endomorphism $N$ with $N^2=0$ on a $d$-dimensional vector
space satisfies
\begin{equation}\label{ir:eq:nilpotent-dimension}
 \dim(\ker N/\operatorname{im} N)=d-2\rank N,
 \qquad
 \dim\operatorname{coker} N=\frac{d+\dim(\ker N/\operatorname{im} N)}2.
\end{equation}

\begin{theorem}[All-field reflection dimension]
\label{ir:thm:field-locus}
Let $q>2$, let $k$ be a field, and use arbitrary weights $a_p\in k$.
If $\operatorname{char}k\ne2$, both sign coinvariants have dimension
$n(q)$. If $\operatorname{char}k=2$, the signs coincide and
\begin{equation}\label{ir:eq:field-locus}
 \dim_k\mathcal D_q(k,a)/(c-1)\mathcal D_q(k,a)
 =n(q)+h(q)\mathbf 1_{\mathcal E_q(a)},
\end{equation}
where $\mathcal E_q(a)$ means
\[
 a_p=1\quad\text{for every odd }p\mid q,
 \qquad a_2\in\{0,1\}\quad\text{if }4\mid q.
\]
There is no condition on $a_2$ when $v_2(q)=1$.
\end{theorem}
\begin{proof}
For characteristic different from two use
Lemma~\ref{ir:lem:balanced}. In characteristic two use the Koszul
calculation and \eqref{ir:eq:nilpotent-dimension} with
$d=\varphi(q)$. Over a field, $a_2(a_2-1)=0$ is equivalent to
$a_2=0$ or $1$.
\end{proof}

Over $\mathbb F_2$, the condition at $2$ is automatic. Over $\mathbb F_4$ it is
not: a weight satisfying $a_2^2+a_2+1=0$ is neither zero nor one,
and the reflection defect disappears when $4\mid q$. Thus an
experiment restricted to binary weights would miss part of the
geometry. The exact regression suite includes all four field
weights at selected levels to test this distinction.

\subsubsection{A complete one-parameter Smith law}
Let $k$ have characteristic two, let $O=k[[t]]$, and choose power
series weights $a_p(t)$. The module $\mathcal D_q(O,a)$ is free of rank
$2n=\varphi(q)$. Form the active sequence $\boldsymbol\delta(t)$
as before, and define its contact order by
\begin{equation}\label{ir:eq:contact-order}
 d=\min_i v_t(\delta_i(t)),\qquad v_t(0)=\infty.
\end{equation}
The word ``contact'' refers to the pullback of the explicit ideal
$I_q$ along the chosen formal curve. It is not an approximation to a
root of a transcendental function.

Each active prime contributes a factor at least two to $\varphi(q)$,
so $\varphi(q)\ge2^{r(q)}$ and $n\ge h$. Thus all block sizes in
the following statement are nonnegative.

\begin{theorem}[Contact-order Smith law]\label{ir:thm:smith-jets}
Let $q>2$, $n=n(q)$, $h=h(q)$, and use the above characteristic-two
power series weights. If $d<\infty$, the Smith diagonal of
$N=c-1$ over $O$ is, up to units,
\begin{equation}\label{ir:eq:smith-contact}
 \operatorname{diag}\bigl(
    \underbrace{1,\ldots,1}_{n-h},
    \underbrace{t^d,\ldots,t^d}_{h},
    \underbrace{0,\ldots,0}_{n}\bigr).
\end{equation}
When $d=0$, the first two blocks together are $n$ unit entries.
When $d=\infty$, replace the $h$ entries $t^d$ by zero; there are
then $n-h$ unit entries and $n+h$ zeros.

For every integer $L\ge1$, put $O_L=k[t]/(t^L)$. The complete jet
cokernel formula is
\begin{equation}\label{ir:eq:jet-length}
 \dim_k\frac{\mathcal D_q(O_L,a)}{(c-1)\mathcal D_q(O_L,a)}
 =nL+h\min(d,L),
\end{equation}
with $\min(\infty,L)=L$.
\end{theorem}

\begin{proof}
First work over $O_L$ and put $d_L=\min(d,L)$. Elementary invertible
operations on the active sequence identify its Koszul complex with
that on $(t^{d_L},0,\ldots,0)$, with the convention that $t^L=0$.
To see this explicitly when $d_L<L$, choose an entry with least
valuation, multiply it by a unit to make it $t^{d_L}$, and subtract
its multiples from all other entries. Every other entry is divisible
by it in $O_L$. If $d_L=L$, all entries already vanish. Invertible
changes of the sequence induce invertible changes on its exterior
algebra, hence isomorphisms of Koszul complexes.

The kernel and cokernel of multiplication by $t^{d_L}$ on $O_L$
each have $k$-dimension $d_L$. Tensoring with $r-1$ zero-differential
two-term complexes multiplies total homology dimension by $2^{r-1}$.
The right-hand side of \eqref{ir:eq:char2-homology} therefore has
$k$-dimension $2^r d_L$. The underlying vector space of
$\mathcal D_q(O_L,a)$ has dimension $2nL$, by the integral basis theorem.
Applying \eqref{ir:eq:nilpotent-dimension} proves
\eqref{ir:eq:jet-length}.

It remains to show that these length formulas determine all Smith
exponents, rather than just their sum. A matrix over the discrete
valuation ring $O$ has nonzero invariant factors $t^{e_i}$ up to
units, together with some zero factors. If there are $z$ zero
factors, its cokernel modulo $t^L$ has length
\[
 zL+\sum_i\min(e_i,L).
\]
The first difference in $L$ is $z+\#\{i:e_i\ge L\}$ for $L\ge1$.
For finite $d$, the formula already proved has first difference
$n+h$ for $1\le L\le d$, and $n$ for $L>d$. Thus exactly $n$
factors are zero and exactly $h$ nonunit factors have exponent $d$;
the remaining $n-h$ factors are units. For $d=0$ there are no
positive exponents. For $d=\infty$, the first difference is always
$n+h$, so there are $n+h$ zero factors and no positive finite
exponents. This proves every case of \eqref{ir:eq:smith-contact}.
\end{proof}

\subsubsection{Examples of the jet law}
At level $15$, take
\[
 a_3(t)=1+t^2,\qquad a_5(t)=1+t^5.
\]
Here $(n,r,h)=(4,2,2)$ and $d=2$. The reflection matrix has Smith
entries $1,1,t^2,t^2,0,0,0,0$, and its jet cokernel dimensions at
$L=1,2,3,4$ are $6,12,16,20$. A rank calculation at $t=0$ alone
would detect the first number, but not the exponent two.

At level $12$, take $a_2(t)=t^3$ and $a_3(t)=1+t^5$.
Then $a_2(a_2-1)=t^3(1+t^3)$, so $d=3$. Since $(n,h)=(2,2)$,
there are no unit Smith factors and the diagonal is
$t^3,t^3,0,0$. Every jet of length at most three has trivial
reflection matrix up to equivalence; later lengths detect the
nonzero factors.

If the entire curve is contained in the exceptional locus, then
$d=\infty$ and the dimension is $(n+h)L$. If any active scalar
has a nonzero constant term, $d=0$ and the dimension is $nL$.
These are exact boundary cases of the same formula. In contrast,
the \emph{unreflected} distribution module always has dimension
$\varphi(q)L$ at length $L$ and never exhibits this jump.

\subsection{Integral jets and why abelian freeness is insufficient}
For a finite free $\Z$-algebra $B$, put
$C_\pm(\mathcal D_B)=\mathcal D_q(B,a)/(c\mp1)\mathcal D_q(B,a)$.
The scalar Tate calculation extends to this coefficient ring.
\begin{proposition}\label{irjet:thm:finite-free-koszul}
With $\overline B=B/2B$ and the reduced active sequence
$\overline{\boldsymbol\delta}$, the two integer-torsion submodules are
\[
 \operatorname{Tor}_{\Z}C_+(\mathcal D_B)
 \simeq\bigoplus_{j\text{ odd}} H_j(K_{\overline B}(\overline{\boldsymbol\delta})),
 \qquad
 \operatorname{Tor}_{\Z}C_-(\mathcal D_B)
 \simeq\bigoplus_{j\text{ even}} H_j(K_{\overline B}(\overline{\boldsymbol\delta})).
\]
These are isomorphisms of $B$-modules. Every torsion class is killed by two.
For $q>2$ and $T=\operatorname{rank}_{\Z}B$, the free abelian rank of
both quotients is $T\varphi(q)/2$.
\end{proposition}
\begin{proof}
The underlying distribution lattice is free over $\Z$ by the all-ring
basis theorem. Thus the torsion lemma still applies. In the fixed-point
periodic complex, multiplication by two is injective on $B$ and its
free chain modules. Its adjacent two-term blocks resolve reduction modulo
two, exactly as in \eqref{ir:eq:integer-tate-koszul}; the bounded chain
width gives the same parity sums as $B$-modules. Over $B\otimes\Q$,
reflection eigenspaces are exact direct summands. The fixed-point Euler
count in the balanced-rank proof is multiplied by $T$ and remains zero,
giving the asserted free ranks. This argument does not assume that the
reflected torsion-free part is a free $B$-module.
\end{proof}

\subsection{An exact Smith formula for arbitrary finite integral jets}
\label{irjet:sec:jets}
Let
\[
 B_L=\Z[u]/(u^L),\quad L\ge1,\qquad
 a_p(u)\in B_L,
\]
and reduce coefficients modulo two. In
$\overline B_L=\mathbb F_2[u]/(u^L)$, form the parameters
\eqref{ir:eq:active-koszul}. For an element $f$ of this ring, define
\[
 \nu(f)=\min\{j:[u^j]f\ne0\},\qquad \nu(0)=L,
 \quad v=\min_{f\in\boldsymbol f}\nu(f).
\]
Here $\boldsymbol f=\overline{\boldsymbol\delta}$ and $r(q)$ is the active-prime count in \eqref{ir:eq:r-h}. For $q>2$, the list is nonempty.

\begin{theorem}[Minimum-valuation jet formula]
\label{irjet:thm:jets}
For $q>2$, the two reflected jet quotients have the following Smith
decomposition as abelian groups:
\begin{equation}
 C_\pm(\mathcal D_{B_L})\simeq
 \Z^{L\varphi(q)/2}\oplus(\Z/2\Z)^{v\,2^{r(q)-1}}.
 \label{irjet:eq:jet-Smith}
\end{equation}
More precisely, their integer-torsion submodules are each isomorphic,
as $B_L$-modules, to
\begin{equation}
 \bigl(B_L/(2,u^v)\bigr)^{2^{r(q)-1}}.
 \label{irjet:eq:jet-torsion-module}
\end{equation}
The direct-sum decomposition in \eqref{irjet:eq:jet-Smith} is a statement over
$\Z$, not a claim that the torsion-free part is free over $B_L$.
\end{theorem}
\begin{proof}
In the principal ideal ring $\overline B_L$, choose a parameter of least
valuation. It is a unit times $u^v$ unless every parameter is zero,
which is the same argument with $v=L$. A change of basis in the free
module defining the Koszul complex transforms the parameter list into
$(u^v,0,\ldots,0)$: first divide by its unit factor, then subtract its
multiples from the other entries. All these are invertible operations
over $\overline B_L$.

The complex is now the tensor product of the two-term complex given by
multiplication by $u^v$ and $r(q)-1$ zero-differential exterior factors.
Both the kernel and the cokernel of multiplication by $u^v$ are
isomorphic to $\overline B_L/(u^v)$; the kernel is the ideal
$u^{L-v}\overline B_L$. Consequently
\[
 H_k\bigl(K_{\overline B_L}(\boldsymbol f)\bigr)
 \simeq\bigl(\overline B_L/(u^v)\bigr)^{\binom{r(q)}{k}}.
\]
Summing either degree parity in Theorem~\ref{irjet:thm:finite-free-koszul} proves
\eqref{irjet:eq:jet-torsion-module}. Each copy has $\mathbb F_2$-dimension $v$,
so Lemma~\ref{ir:lem:torsion-tate} gives the torsion multiplicity in
\eqref{irjet:eq:jet-Smith}.

The unreflected underlying lattice has rank $L\varphi(q)$.
The alternating fixed-point trace used in
Theorem~\ref{ir:thm:integer-reflection} is multiplied by $L$, and still
vanishes. Thus each rational eigenspace has dimension
$L\varphi(q)/2$. This proves the free rank and completes the Smith
calculation.
\end{proof}

\begin{example}[Five jet coefficients at level fifteen]
Take $L=5$, $a_3(u)=1+u^2$, and $a_5(u)=1+2u+u^3$.
The binary parameters are $u^2,u^3$, so $v=2$ and $r(q)=2$. Therefore
\[
 C_\pm(\mathcal D_{B_5})\simeq\Z^{20}\oplus(\Z/2\Z)^4.
\]
The constant jets $a_3=a_5=1$ would instead give ten two-torsion factors.
The constant terms of the weights agree; their higher binary
coefficients change the integral reflection structure.
\end{example}

\begin{example}[The two-adic jet parameter matters]
At $q=12$, $L=4$, let $a_2(u)=1+u$ and $a_3(u)=1+u^3$.
The two parameters are $u^3$ and $u+u^2$, not simply
$1-a_3$ and $1-a_2$. Thus $v=1$, $r(q)=2$, and
\[
 C_\pm(\mathcal D_{B_4})\simeq\Z^8\oplus(\Z/2\Z)^2.
\]
Although the parity of a \emph{scalar} two-prime weight does not affect
Theorem~\ref{ir:thm:integer-reflection}, a nonconstant two-prime jet can affect
Theorem~\ref{irjet:thm:jets}.
\end{example}

\subsubsection{Why the free abelian part need not be jet-free}
Take $q=3$, $L=2$, and $a_3=1+u$. In the unreflected basis
$e=e_0$, $f=e_{2/3}$, reflection satisfies
\[
 ce=e,\qquad cf=ue-f.
\]
Hence
\begin{equation}
 C_+(\mathcal D_{B_2})=B_2e\oplus B_2f\big/\langle ue-2f\rangle.
 \label{irjet:eq:nonflat}
\end{equation}
As an abelian group this has free generators $e,f$ and a torsion generator
$uf$ of order two. In its torsion-free quotient, multiplication by $u$
sends $e$ to $2f$ and $f$ to zero. This is not a free rank-one
$B_2$-module: the image of $u$ in a free rank-one module is a primitive
rank-one sublattice, whereas here the image is $2\Z f$.
Thus a claim of $B_L$-freeness for the reflected torsion-free part would
be false. By contrast, $C_-=B_2f\oplus(B_2/(2,u))e$ in this example.
Both signs have the same abelian Smith factors, but different module
extensions over the jet ring.

\subsubsection{Multivariate jets and arithmetic weights}
For a finite free $\Z$-algebra $B$ of rank $L$, the Koszul theorem remains
an exact answer. If $q>2$ and
$h_k=\dim_{\mathbb F_2} H_k(K_{B/2B}(\boldsymbol f))$, then
\[
 \operatorname{Tor}_\Z C_+\simeq(\Z/2\Z)^{\sum_{k\text{ odd}}h_k},\qquad
 \operatorname{Tor}_\Z C_-\simeq(\Z/2\Z)^{\sum_{k\text{ even}}h_k}.
\]
The two sums are equal: the Euler characteristic of the Koszul complex
is $L(1-1)^{r(q)}=0$. A closed minimum-valuation description is special
to the one-variable principal ideal ring; it must not be transferred
unchanged to a multivariate Artinian ring.

For actual positive integral polylogarithm orders, $A_p=p^{1-n}$ usually
belongs to $\Z[1/q]$, rather than $\Z$. The corresponding coefficient
ring must be stated. If $q$ is even, inverting $q$ also inverts two and
removes the reflected two-torsion. For odd $q$, every $p^{1-n}$ is one
modulo two, and the ordinary maximal multiplicity persists over
$\Z[1/q]$. These are statements about the chosen formal coefficient
lattice, not about independence of the evaluated periods.

\subsection{Separable multivariable jets: an exact product law}
\label{irmulti:sec:product}
The one-variable minimum-valuation formula has a useful multivariable
extension when the active parameters separate into coordinate powers.
The general Koszul answer remains necessary for parameters that do not
have this form. In particular, replacing several coordinate valuations
by their minimum would give the wrong answer here.

Let $q>2$, put $r=r(q)$, $n=\varphi(q)/2$, $h=2^{r-1}$, and set
\[
 S=k[t_1,\ldots,t_r]/(t_1^{L_1},\ldots,t_r^{L_r}),\qquad
 L_i\ge1,\qquad \operatorname{char}k=2.
\]
Suppose the active sequence, after permutation and multiplication of
individual entries by units, is $(t_1^{d_1},\ldots,t_r^{d_r})$,
where $d_i\ge1$; a zero entry is allowed by taking $d_i=\infty$.
Write $v_i=\min(d_i,L_i)$ and $T=\prod_iL_i$.

\begin{theorem}[Multivariable reflection product law]
\label{irmulti:thm:dimension}
Under these hypotheses,
\begin{align}
 \dim_k H_j(K_S(\boldsymbol\delta))
   &=\binom rj\prod_{i=1}^r v_i,\qquad 0\le j\le r,
       \label{irmulti:eq:homology}\\
 \dim_k\mathcal D_q(S,a)/(c-1)\mathcal D_q(S,a)
   &=nT+h\prod_{i=1}^r v_i.\label{irmulti:eq:dimension}
\end{align}
The unreflected dimension is $\varphi(q)T$.
For integral weights in
$B=\Z[t_1,\ldots,t_r]/(t_1^{L_1},\ldots,t_r^{L_r})$
whose reduced active sequence has the above form, both reflected
coinvariants have the abelian-group decomposition
\[
 C_\pm(\mathcal D_B)\simeq
 \Z^{nT}\oplus(\Z/2\Z)^{h\prod_i v_i}.
\]
Their integer-torsion submodules are, more precisely, isomorphic as
$B$-modules to
\[
 \left(B/(2,t_1^{v_1},\ldots,t_r^{v_r})\right)^h.
\]
The integral assertion uses $k=\mathbb F_2$.
\end{theorem}
\begin{proof}
Multiplying an active entry by a unit and permuting entries induces an
invertible change of the exterior generators of its Koszul complex.
After that change, the complex is the tensor product over $k$ of the
two-term complexes
\[
 k[t_i]/(t_i^{L_i})\xrightarrow{t_i^{d_i}}k[t_i]/(t_i^{L_i}).
\]
Each has kernel and cokernel of dimension $v_i$, in degrees one and
zero respectively. As a module over its truncated polynomial ring,
each homology group is isomorphic to $k[t_i]/(t_i^{v_i})$; multiplication
by $t_i^{L_i-v_i}$ gives the kernel isomorphism. This includes the
zero-map case $v_i=L_i$.

The K\"unneth formula over the field $k$ has no Tor corrections. In
total degree $j$ there is one tensor summand for each $j$-element subset
of the $r$ factors. Each summand has dimension $\prod_i v_i$ and is
isomorphic, as an $S$-module, to
$S/(t_1^{v_1},\ldots,t_r^{v_r})$.
This proves \eqref{irmulti:eq:homology} and its module refinement.
The characteristic-two fixed-point theorem gives total reflection
homology of dimension $2^r\prod_i v_i$. Since $\mathcal D_q(S,a)$
has dimension $2nT$, the square-zero rank identity
\eqref{ir:eq:nilpotent-dimension} gives \eqref{irmulti:eq:dimension}.
For the integral statement apply the finite-free-algebra Tate formula
of Proposition~\ref{irjet:thm:finite-free-koszul}. The even and odd
binomial sums are both $2^{r-1}$; the balanced-rank calculation gives
free rank $nT$. Finally a finitely generated abelian group splits into
its free part and its torsion subgroup. No splitting as $B$-modules
or freeness of its torsion-free part is claimed.
\end{proof}

\begin{example}[Two independent deformation directions at level fifteen]
Take $B=\Z[t,u]/(t^3,u^4)$,
$a_3=1+t^2$ and $a_5=1+u^3$. Then $n=4$, $r=2$, $h=2$,
$T=12$ and $(v_1,v_2)=(2,3)$. Both integral signs give
\[
 C_\pm(\mathcal D_B)\simeq\Z^{48}\oplus(\Z/2\Z)^{12},
\]
and the characteristic-two reflected dimension is $60$.
The defect depends on the product $2\cdot3$, not on a single minimum.
At an active two-prime coordinate one may take $a_2=t_i^{d_i}$
or $a_2=1+t_i^{d_i}$: in either case
$a_2(a_2-1)=t_i^{d_i}(1+t_i^{d_i})$, and the second factor is a unit.
\end{example}

The independent raw-matrix verifier expands the original prime rows and
reflection rows in point/monomial coordinates over $\mathbb F_2$.
It imports neither the normal-form generator nor the Koszul calculation.
It checks 210 finite cases, including both two-prime branches, vanished
coordinate powers, unequal jet lengths, and a three-prime level. A
corruption control rejects omission of reflection rows. These computations
corroborate the formula; the proof above establishes its full scope.

\subsection{Exact cyclotomic identities and pole-cancelled spectral jets}
\label{ir:sec:identities}

\subsubsection{Evaluating an integral certificate}
For $x\in\Q/\Z$, write
\[
 \intdistL{s}{x}=\Li_s(e^{2\pi i x}).
\]
At $x=0$, this means $\zeta(s)$. For $x\ne0$, the spectral function
is entire. One way to see this, and to fix its continuation, is the
finite residue-class formula
\begin{equation}\label{ir:eq:fourier-hurwitz}
 \intdistL{s}{a/q}=q^{-s}\sum_{u=1}^{q}e^{2\pi i au/q}\zeta(s,u/q).
\end{equation}
For $a\not\equiv0\pmod q$, the residues of the simple poles at
$s=1$ sum to zero. The standard meromorphic continuation of Hurwitz
zeta has no other spectral poles \cite{intdist:DLMFHurwitz}.

For $\Re s>1$, absolute convergence gives
\begin{equation}\label{ir:eq:analytic-distribution}
 \sum_{py=x}\intdistL{s}{y}=p^{1-s}\intdistL{s}{x}.
\end{equation}
Indeed the root sum cancels every series index not divisible by $p$;
for index $pk$ it contributes $p e^{2\pi i kx}$.
Continuation then proves the meromorphic identity for every complex
$s$. Thus an integral polynomial certificate becomes an exact analytic
identity under
\begin{equation}\label{ir:eq:polylog-weights}
 e_x\longmapsto\intdistL{s}{x},\qquad A_p\longmapsto p^{1-s},
\end{equation}
where real positive $p$ fixes the exponential unambiguously. No integer
relation search or floating-point acceptance threshold occurs in this
step.

\begin{proposition}[A removable endpoint coefficient]
\label{ir:prop:endpoint-zero}
For $y\ne0$, the coefficient $c_{y,0}(\mathbf A)$ in its integral
normal form vanishes at $A_p=1$ for all $p$. Hence its polylogarithmic
specialization cancels the simple pole of $\zeta(s)$ at $s=1$.
\end{proposition}
\begin{proof}
At ordinary weights, the functional taking $e_0$ to $1$ and all
other grid symbols to $0$ annihilates every distribution row. For
$x=0$, both the root sum and $e_x$ contain exactly one copy of
$e_0$; for $x\ne0$, neither does. Apply the functional to the
normal form of $e_y$. The second assertion follows because
$p^{1-s}=1$ at $s=1$ and the coefficient is analytic in $s$.
\end{proof}
The vanishing need only be first order. It is not legitimate to infer
one vanishing factor for each prime merely from the number of prime
variables in a raw-point normal form.

\subsubsection{A three-row level-twelve identity}
In the level-$12$ free module, define
\[
 F_{12}=e_{1/12}+e_{5/12}+(A_3+1)e_{3/4}
        -A_2A_3(A_2-1)e_0.
\]
The following is an equality before taking the quotient:
\begin{equation}\label{ir:eq:cert12}
 F_{12}=r_{3,1/4}+A_3r_{2,1/2}+A_2A_3r_{2,0}.
\end{equation}
For example, the first row contains
$e_{1/12}+e_{5/12}+e_{3/4}-A_3e_{1/4}$; the next two rows cancel
$e_{1/4}$ and then $e_{1/2}$. This verifies the certificate by direct
inspection.

\begin{corollary}\label{ir:cor:level12}
For every complex $s$, with the removable value used at $s=1$,
\begin{align}
 \intdistL{s}{1/12}+\intdistL{s}{5/12}
 +(3^{1-s}+1)\intdistL{s}{3/4}
 =6^{1-s}(2^{1-s}-1)\zeta(s).
 \label{ir:eq:level12}
\end{align}
At $s=1$, the right-hand side is $-\log2$.
\end{corollary}
\begin{proof}
Evaluate \eqref{ir:eq:cert12} using
\eqref{ir:eq:polylog-weights}. The value at one follows from
$2^{1-s}-1=-(s-1)\log2+O((s-1)^2)$.
\end{proof}

\subsubsection{A four-row level-fifteen identity}
Set
\begin{equation}\label{ir:eq:U-set}
 U=\{1/15,4/15,7/15,2/5,4/5\}.
\end{equation}
In the level-$15$ free module define
\begin{align}
 F_{15}={}&e_{8/15}-\sum_{u\in U}e_u
 -(A_3+1)e_{3/5}-(A_5+1)e_{2/3}
 -(1-A_3A_5)e_0.
 \label{ir:eq:F15}
\end{align}
A short exact certificate is
\begin{equation}\label{ir:eq:cert15}
 F_{15}=r_{3,3/5}-r_{5,1/3}-A_5r_{3,0}-r_{5,0}.
\end{equation}
The roots in the first two rows are respectively
\[
 \{1/5,8/15,13/15\},\qquad
 \{1/15,4/15,7/15,2/3,13/15\}.
\]
The $13/15$ terms cancel. The last two rows remove $1/3$ and
$1/5$, and leave precisely \eqref{ir:eq:F15}. In particular the
coefficient of $e_0$ in the normal form of $e_{8/15}$ is
$1-A_3A_5$, not $(A_3-1)(A_5-1)$.

\begin{theorem}[All-order level-fifteen identity]
\label{ir:thm:level15}
For every complex spectral order $s$,
\begin{align}
 \intdistL{s}{8/15}-\sum_{u\in U}\intdistL{s}{u}
 &-(3^{1-s}+1)\intdistL{s}{3/5}
  -(5^{1-s}+1)\intdistL{s}{2/3}\nonumber\\
 &=(1-15^{1-s})\zeta(s).
 \label{ir:eq:level15}
\end{align}
Both sides extend holomorphically through $s=1$, where the right-hand
side is $\log15$.
\end{theorem}
\begin{proof}
Apply \eqref{ir:eq:polylog-weights} to the four-row identity.
The endpoint value follows from
$1-15^{1-s}=(s-1)\log15+O((s-1)^2)$ and the residue one of
$\zeta(s)$.
\end{proof}
The theorem supplies infinitely many exact identities, including
ordinary positive integer weights, negative orders, and derivatives
in spectral order. It is generated by classical distribution rather
than by a new numerical relation conjecture.

A compatible level-$30$ normal form follows from the extra row
$r_{2,1/15}$:
\begin{align}
 e_{1/30}={}&(A_2-1)e_{1/15}-(A_3+1)e_{3/5}
 -(A_5+1)e_{2/3}\nonumber\\
 &-e_{4/15}-e_{7/15}-e_{2/5}-e_{4/5}
 +(A_3A_5-1)e_0.\label{ir:eq:level30-normal}
\end{align}
Its residual has the five-row certificate
\[
 r_{2,1/15}-r_{3,3/5}+r_{5,1/3}+A_5r_{3,0}+r_{5,0}.
\]
All rows can be read in the level-$30$ grid. This illustrates
compatibility of the straightener with an enlarged level and a
weight whose order-one specialization vanishes in a particular
coefficient.

\subsubsection{Every spectral derivative, with the pole term retained}
Let
\[
 \intdistJ{j}{x}=\left.\frac{\partial^j}{\partial s^j}\intdistL{s}{x}
                   \right|_{s=1},\qquad x\ne0,
\]
and write the Laurent expansion
\begin{equation}\label{ir:eq:stieltjes}
 \zeta(1+t)=\frac1t+\sum_{m\ge0}\frac{(-1)^m\gamma_m}{m!}t^m.
\end{equation}
For a positive prime $p$ and $N\ge0$, define the explicit finite operator
\begin{equation}\label{ir:eq:T-operator}
 \mathcal T_{p,N}(x)=\intdistJ{N}{x}
 +\sum_{j=0}^{N}\binom Nj(-\log p)^j\intdistJ{N-j}{x}.
\end{equation}
It is the $N$th derivative of
$(p^{1-s}+1)\intdistL{s}{x}$ at $s=1$.

\begin{theorem}[A level-fifteen Stieltjes--polylogarithm jet family]
\label{ir:thm:jets15}
For every integer $N\ge0$, put $\ell=\log15$. Then
\begin{align}
 &\intdistJ{N}{8/15}-\sum_{u\in U}\intdistJ{N}{u}
 -\mathcal T_{3,N}(3/5)-\mathcal T_{5,N}(2/3)\nonumber\\
 &\qquad =(-1)^N\left\{\frac{\ell^{N+1}}{N+1}
 -\sum_{j=1}^{N}\binom Nj\ell^j\gamma_{N-j}\right\}.
 \label{ir:eq:jets15}
\end{align}
The sum on the right is empty at $N=0$.
\end{theorem}
\begin{proof}
All nonzero-root terms of \eqref{ir:eq:level15} are entire in $s$,
so differentiate the holomorphic continuation $N$ times. The weighted
terms give \eqref{ir:eq:T-operator} by Leibniz's rule. For the right
side use
\[
 1-e^{-\ell t}=\sum_{j\ge1}\frac{(-1)^{j+1}\ell^j}{j!}t^j
\]
and multiply by \eqref{ir:eq:stieltjes}. The pole term contributes
$(-1)^N\ell^{N+1}/(N+1)$ to the $N$th derivative. The regular
Laurent coefficients contribute
$(-1)^{N+1}\sum_{j=1}^N\binom Nj\ell^j\gamma_{N-j}$.
This is the stated formula.
\end{proof}

The first derivative gives the especially compact identity
\begin{align}
 \intdistJ{1}{8/15}-\sum_{u\in U}\intdistJ{1}{u}
 -2\intdistJ{1}{3/5}-2\intdistJ{1}{2/3}
 &+\log3\,\intdistJ{0}{3/5}+\log5\,\intdistJ{0}{2/3}\nonumber\\
 &=\gamma_0\log15-\tfrac12\log^2 15.
 \label{ir:eq:first-jet}
\end{align}
For $N=2$, the scalar right-hand side is
\[
 \frac{\log^3 15}{3}-\gamma_0\log^2 15-2\gamma_1\log15.
\]
These terms would be wrong if one substituted $s=1$ into the
coefficient of $\zeta(s)$ before cancelling its pole.

More generally, if an analytic coefficient $c(t)$ satisfies $c(0)=0$,
then
\begin{equation}\label{ir:eq:general-pole-jet}
 \left.\frac{d^N}{dt^N}\{c(t)\zeta(1+t)\}\right|_{t=0}
 =\frac{c^{(N+1)}(0)}{N+1}
 +\sum_{j=0}^{N}\binom Nj c^{(j)}(0)(-1)^{N-j}\gamma_{N-j}.
\end{equation}
This coefficient identity gives a general implementation rule for every
raw-point certificate, not just the displayed level.

\subsubsection{An elementary endpoint and branch check}
For $0<x<1$, the chosen branch gives
\[
 \intdistJ{0}{x}=-\log(2\sin\pi x)+i\pi(1/2-x).
\]
Taking real parts of \eqref{ir:eq:level15} at $s=1$ gives
\[
 \frac{\prod_{u\in U}(2\sin\pi u)
 (2\sin(3\pi/5))^2(2\sin(2\pi/3))^2}
 {2\sin(8\pi/15)}=15.
\]
Since $\sin(8\pi/15)=\sin(7\pi/15)$, it simplifies to
\begin{equation}\label{ir:eq:sine-product}
 (2\sin(\pi/15))(2\sin(4\pi/15))
 (2\sin(2\pi/5))^3(2\sin(\pi/5))=5.
\end{equation}
All factors are positive, so exponentiating the real logarithmic identity
introduces no ambiguity. The imaginary parts cancel by the displayed
linear argument formula. This is an elementary specialization and branch
check, not a claim of novelty for the trigonometric product itself.

\subsubsection{A division-free Hurwitz--Stieltjes translation}
The same polynomial certificates evaluate under
\[
 A_p=p^s,\qquad e_x\mapsto Z_x(s),\qquad
 Z_x(s)=\begin{cases}\zeta(s,x),&0<x<1,\\\zeta(s,1),&x=0.\end{cases}
\]
The distribution identity here is
$\sum_{py=x}Z_y(s)=p^sZ_x(s)$; it follows by splitting the Hurwitz
series into residue classes, including the endpoint convention.
Let $c_b(t)$ be the normal-form coefficient of $e_b$ for $e_x$ after
$A_p=p^{1+t}$, and put $C(t)=\sum_b c_b(t)$.
Residue comparison gives $C(0)=1$. If
$\gamma_j(b^*)$ denotes the generalized Stieltjes coefficient with
$b^*=b$ for $b\ne0$ and $b^*=1$ for $b=0$, expansion of the exact
normal form gives, for all $N\ge0$,
\begin{align}
 \gamma_N(x^*)={}&
 \frac{(-1)^N C^{(N+1)}(0)}{N+1}
 +\sum_{b\in\mathcal B_q}\sum_{j=0}^{N}
 (-1)^j\binom Nj c_b^{(j)}(0)\gamma_{N-j}(b^*).
 \label{ir:eq:hurwitz-jets}
\end{align}
To verify the sign, multiply each Laurent series by $c_b(t)$, compare
$t^N$, and multiply by $(-1)^NN!$. The pole part is $C(t)/t$ and
contributes the first term. Thus even the underived constants require
a coefficient derivative. This is the raw-basis version of the
manuscript's pole-corrected jet philosophy, now with division-free
integral certificates for the coefficients.

Every coefficient derivative in \eqref{ir:eq:hurwitz-jets} is a finite
integer combination of products of logarithms of primes and integer
prime powers. There is no need to differentiate a solved matrix inverse
or to redo a relation search separately at each derivative order.

\subsubsection{Level-twelve derivatives and cyclotomic products}
For the coefficient $c(t)=12^{-t}-6^{-t}$ in
\eqref{ir:eq:level12}, set $D_j=(-\log12)^j-(-\log6)^j$.
The finite-part rule \eqref{ir:eq:general-pole-jet} yields, for every $N\ge0$,
\begin{align}
 \intdistJ{N}{1/12}+\intdistJ{N}{5/12}
 +\mathcal T_{3,N}(3/4)
 &=\frac{D_{N+1}}{N+1}
 +\sum_{j=1}^N\binom Nj D_j(-1)^{N-j}\gamma_{N-j}.
\end{align}
For example its first derivative has right side
$(\log^2 12-\log^2 6)/2-\gamma_0\log2$.
At order one of the polylogarithm, exponentiating the exact logarithmic
identities gives, with $\zeta_q=e^{2\pi i/q}$,
\begin{align}
 (1-\zeta_{12})(1-\zeta_{12}^5)(1-\zeta_{12}^9)^2&=2,\\
 (1-\zeta_{30})(1-\zeta_{30}^8)(1-\zeta_{30}^{12})
 (1-\zeta_{30}^{14})(1-\zeta_{30}^{18})^2
 (1-\zeta_{30}^{20})^2(1-\zeta_{30}^{24})&=15.
\end{align}
The level-thirty identity is the analytic specialization of
\eqref{ir:eq:level30-normal}; its endpoint coefficient has removable
value $-\log15$. Exponentiation proves the product without a branch
ambiguity, while the converse would lose an additive multiple of
$2\pi i$. The preserved independent polynomial verifier also proves
that each product polynomial minus its displayed value is divisible
by the corresponding cyclotomic polynomial.

